\documentclass[11pt]{article}
\usepackage[margin=1in]{geometry}
\usepackage{parskip}
\usepackage{enumitem}
\setlist{nosep}

\title{\vspace{-1.5cm}\textbf{Working with the xArm 6 Robot Arm}}
\author{Project Report}
\date{\today}

\begin{document}
\maketitle
\vspace{-0.8cm}

\section*{What We Set Out to Do}

The goal was to teach a robot arm (the xArm~6) to do useful jobs on its own.
We started simple: pick up a box from one place and set it down in another.
Then we made it smarter: use a camera so the arm can \emph{see} a red cube,
figure out where it is and how big it is, go to it by itself, and pick it up.

\section*{Step One: Pick and Place}

First we taught the arm two spots by hand --- where to grab the box, and where
to put it down. We then wrote a simple program so that, with one click, the arm
would move to the box, close its gripper, carry it across, and release it, then
return to its resting position. To make this easy, everything was set up as
click-to-run buttons on the desktop, so no typing is needed.

\section*{The Problems We Ran Into (and Fixed)}

The arm did not work perfectly at first, and each problem taught us something:

\begin{itemize}
  \item \textbf{The arm kept stopping suddenly.} It turned out the arm did not
        know how heavy its own gripper was, so it \emph{thought} it was bumping
        into something and stopped for safety. Once we told it the gripper's
        weight, this stopped happening.
  \item \textbf{It froze when reaching too far.} Some target spots were near the
        very edge of how far the arm can stretch. We changed the way it moves so
        it takes smoother paths and checks in advance that it can actually reach
        a spot before trying.
  \item \textbf{The gripper sat crooked.} We straightened it and set its
        resting angle to a clean zero.
\end{itemize}

Every version of the program was carefully double-checked for mistakes before
being used, since a real robot is involved and safety comes first.

\section*{Step Two: Giving the Arm Eyes}

Next we added a small depth camera to the arm's wrist so it could find a red
cube by itself. The arm learns the link between what the camera sees and where
things really are by making a couple of small test moves and watching how the
view changes --- it calibrates itself, instead of us measuring by hand. It also
adjusts its grip to match the size of each cube automatically.

To make the fiddly steps easier, we set it up so the arm can be driven around
with the keyboard (like a simple game) instead of typing in numbers.

\section*{Where Things Stand Now}

The camera can reliably spot the red cube and measure it. The arm connects and
moves safely, with several protections built in so it stops cleanly and reports
plainly if anything goes wrong. The final calibration step --- lining the
fingers up around the cube once so the arm learns its reference --- is the last
piece to complete before fully hands-free, automatic picking is ready to run.

\end{document}
